Soit $n$ un entier naturel non nul. Un \emph{carré nordique} est une grille $n \times n$ dont les cases ont
été numérotées de $1$ à $n^2$ dans un ordre arbitraire, chaque numéro étant utilisé une seule fois. Deux
cases distinctes sont considérées comme adjacentes si elles ont un côté commun. Toute case dont le numéro
est inférieur aux numéros des cases qui lui sont adjacentes est appelée \emph{vallée}.
Un \emph{chemin de randonnée} est une suite formée d'au moins une case, telle que :
\begin{enumerate}
  \item[(i)] la première case de la suite est une vallée ;
  \item[(ii)] deux cases consécutives de la suite sont toujours adjacentes ;
  \item[(iii)] les cases de la suite sont numérotées dans l'ordre strictement croissant.
\end{enumerate}
Trouver, en fonction de $n$, le plus petit nombre possible de chemins de randonnée que peut contenir un
carré nordique.
